% Three-page Results revision. Evidence: e267cca82b5104c125dc5e2c8d338c80b486e5de.
% Use with the accompanying appendix, figures and tables.
\section{Results}
\label{sec:results}

We examine inclusion sensitivity, its relationship to effectiveness, and the behavior of order-handling controls. Rates are query-macro percentages on the 308 WANDS and 499 ESCI eligible queries unless stated otherwise; differences are percentage points (pp). nDCG uses the $0$--$1$ scale. Comparisons use paired query-cluster intervals.

\subsection{Attribute-Order Sensitivity (RQ1)}
\label{sec:results_rq1}

Attribute order changes relevant-product inclusion in every dataset--encoder setting. Across the seven catalog-wide schedules, VI@20 ranges from 11.93--19.58\% on WANDS and 3.66--6.91\% on ESCI. VI describes crossing within the tested ordering family, rather than the probability of omission under future orders.

Changing only the target representation is sufficient. With all competitors fixed at source order, VI@20 remains 10.83--19.45\% on WANDS and 2.42--5.88\% on ESCI (Figure~\ref{fig:rq1_target_fitting}). The two interventions answer complementary questions; their difference does not decompose target and competitor effects. Sensitivity also persists when every tested target input fits completely: fully-fitting VI@20 ranges from 7.16--14.55\% and 2.55--5.88\%, respectively. All catalog inputs fit GTE's token limit. The restricted populations differ across encoders, so these comparisons do not identify a causal truncation effect. They establish that truncation is not necessary for inclusion sensitivity.

\begin{figure*}[!tbp]
  \centering
  \includegraphics[width=\textwidth]{figures/rq1_vi_top20_main.pdf}
  \caption{Target-only VI@20 on full and fully-fitting supports, with 95\% intervals and query/pair counts. Competitors remain fixed; panel scales differ.}
  \label{fig:rq1_target_fitting}
  \Description{Three encoders retain nonzero cutoff crossing on both datasets when complete target inputs fit. Supports are reported separately.}
\end{figure*}

Larger ordering families reveal additional crossing. On a fixed 128-query WANDS supplement, expanding from 7 to 32 schedules raises target-only VI@20 from 20.85\% to 28.21\% for MiniLM and 11.40\% to 17.55\% for BGE. Nested VI is nondecreasing by definition; the empirical result is the additional affected support. Appendix~\ref{app:ch6_coverage} reports fitting controls, ESCI target-order enumeration, and execution-related numerical limits. Native score-boundary cases and the limitations of attribution-based explanations appear in Appendices~\ref{app:ch6_boundaries} and~\ref{app:ch6_xai}.

\subsection{Aggregate Effectiveness and Product Membership (RQ2)}
\label{sec:results_rq2}

Recall measures the balance between newly included and newly omitted relevant products; their sum measures membership change (Figure~\ref{fig:rq2_membership}). Reversing WANDS attributes with BGE newly includes 2.357\% and newly omits 2.263\% of relevant support. Their combined share is 4.621\%, although Recall changes by only $+0.094$ pp (95\% CI $[-0.409,0.563]$). Each query's denominator is $|H_q|$. Cancellation also occurs within queries: across the six source-order comparisons, 67--78 of the 308 WANDS queries have equal positive BGE gain and loss counts. Their Recall is exactly unchanged despite different retrieved products. These are ranges across comparisons, not disjoint groups.

\begin{figure*}[!tbp]
  \centering
  \includegraphics[width=\textwidth]{figures/rq2_membership_top20_main.pdf}
  \caption{Membership changes from source order at $K=20$. Bars show newly omitted/included relevant shares; diamonds show net Recall changes with 95\% paired intervals.}
  \label{fig:rq2_membership}
  \Description{Across six dataset--encoder settings, relevant-product gains and losses can offset each other in aggregate Recall.}
\end{figure*}

Graded effectiveness is not generally unchanged. Among queries with changed Top-20 relevant-product membership, median absolute nDCG@20 differences range from 0.0226--0.0399 on WANDS and 0.0388--0.0482 on ESCI. Query-mean differences are smaller: $-0.01101$ to $+0.00151$ and $-0.00393$ to $+0.00092$, respectively. Five of the 36 uncorrected paired intervals are entirely negative; the others contain zero, without establishing equivalence. Aggregate effectiveness therefore complements, but does not fully describe, product membership. Appendix~\ref{app:ch6_graded} reports conditional distributions, numerical equality, condensed cNDCG and judged coverage.

\subsection{Effectiveness and Consistent Inclusion under Controls (RQ3)}
\label{sec:results_rq3}

Table~\ref{tab:strategy_joint} evaluates controls on common populations. Raw dense and raw hybrid effectiveness averages seven schedules. Canonical, set-based, fixed-view and lexical controls produce order-independent indexes, making their zero VI structural. Their effectiveness and the identities of persistently included products remain empirical questions. All three canonical rules are retained without selecting a test-set winner.

\begin{table*}[t]
\centering
\small
\setlength{\tabcolsep}{4pt}
\caption{Joint strategy evaluation: mean Recall@100 / VI@100 (\%).}
\label{tab:strategy_joint}
\begin{tabular}{lrrrr}
\toprule
 & \multicolumn{2}{c}{WANDS} & \multicolumn{2}{c}{ESCI} \\
Method & MiniLM & BGE & MiniLM & BGE \\
\midrule
Raw (seven schedules) & 56.61 / 21.89 & 63.62 / 12.30 & 86.19 / 1.41 & 91.04 / 0.72 \\
Canonical ascending & 56.05 / 0.00 & 63.24 / 0.00 & 86.41 / 0.00 & 91.05 / 0.00 \\
Canonical descending & 56.16 / 0.00 & 63.61 / 0.00 & 86.38 / 0.00 & 91.01 / 0.00 \\
Canonical type/style & 57.19 / 0.00 & 63.56 / 0.00 & 86.36 / 0.00 & 91.09 / 0.00 \\
Set-Mean & 58.81 / 0.00 & 61.82 / 0.00 & 84.50 / 0.00 & 88.95 / 0.00 \\
Centroid, $M=2$ & 56.44 / 0.00 & 63.96 / 0.00 & 86.33 / 0.00 & 91.07 / 0.00 \\
Multi-vector max, $M=2$ & 57.04 / 0.00 & 64.00 / 0.00 & 86.35 / 0.00 & 91.09 / 0.00 \\
BM25 & 63.23 / 0.00 & 63.23 / 0.00 & 82.83 / 0.00 & 82.83 / 0.00 \\
Hybrid (seven schedules) & 65.55 / 11.08 & 67.77 / 7.05 & 89.57 / 0.56 & 91.61 / 0.74 \\
\bottomrule
\end{tabular}
\par\smallskip\begin{minipage}{\linewidth}\small Raw and hybrid average all seven schedules. Canonical, Set-Mean, canonical-seeded centroid/max and BM25 use a fixed input-order independent index: their per-schedule Recall is identical and structural VI is zero. All three canonical rules are retained. WANDS: 308 queries, 21,299 Exact pairs; ESCI: 499 queries, 4,434 E pairs. Source-order Recall, persistent states, paired intervals, membership counts and cost are reported in the artifact tables.\end{minipage}
\end{table*}


Zero VI can fix an unstable product to either outcome (Table~\ref{tab:control_destinations}). Under WANDS/MiniLM ascending canonicalization at $K=20$, 7.57\% of all relevant support moves from crossing to fixed inclusion, while 12.01\% moves to fixed omission. The control additionally omits 0.23\% that was included under every raw schedule. These are query-macro shares of all relevant support, not conditional failure rates. Canonicalization removes order dependence without necessarily preserving relevant coverage; complete transitions appear in Appendix~\ref{app:ch6_transitions}.

\input{tables/control_destinations_main.tex}

Effectiveness varies across settings. Set-Mean exceeds raw dense Recall@100 on WANDS/MiniLM but falls below it in the other three settings. Raw hybrid improves Recall@100 while reducing VI@100 on WANDS/MiniLM, WANDS/BGE and ESCI/MiniLM, with paired intervals supporting both directions. ESCI/BGE differences remain inconclusive. Raw hybrid retains nonzero VI in all four settings.

Canonical--BM25 hybrids combine lexical evidence with order invariance. All 12 dataset--encoder--rule configurations have positive uncorrected nDCG@20 intervals versus BM25, and nine do versus their corresponding canonical dense branch. Against the raw-hybrid seven-schedule mean, two nDCG@20 contrasts are positive and ten are inconclusive; all Recall@100 intervals include zero. This does not establish equivalence or universal preservation: WANDS/BGE ascending reduces Recall@20 by 0.624 pp (95\% CI $[-1.370,-0.084]$). The results support competitive simple controls, rather than a general effectiveness--instability trade-off.

Cost further distinguishes the controls. Canonical and centroid retrieval retain one dense vector per product; multi-vector retrieval stores and scores $M$ vectors, and hybrid retrieval adds a sparse index and branch scoring. Appendix~\ref{app:ch6_transitions} provides view-budget and timing results. The practical comparison therefore requires effectiveness, persistent inclusion and omission, identity transitions, and cost alongside VI.
